\scanpage{103}
\mspage{96}

Mais revenons à ce cas

\[
 X=\varinjlim X_i\ \bigg|\ \begin{aligned}&X_i\in M^\pi\\&I\text{ grand devant }\pi,\end{aligned}
\]

et introduisons à nouveau

\[
 \widetilde I=\{I'\subset I\mid I'\text{ filtrant, de cardinal }\leq\pi\},
\]

\unsure{donc|alors} cette fois $\widetilde I$ est \deleted{\unsure{pr}} bien adapté à $\pi$ (filtrant, et toute partie \inserted{filtrante} de cardinal $\leq\pi$ a une borne supérieure), et si on pose

\[
 X_{\widetilde i}=\underset{i\in I'_{\widetilde i}}{\varinjlim} X_i,
\]

on voit que ce syst. inductif

\[
 (X_{\widetilde i})_{\widetilde i\in\widetilde I}
\]

est \unsure{parfait\textsuperscript{mt}} adapté à $\pi$ et la filtr. cardinale. De plus, \deleted{\unsure{on a} une application canonique}

\deleted{\[
 I\longrightarrow\widetilde I\qquad i\longmapsto\{i\}
\]}

\deleted{(pas croissante, qu'à cela ne tienne), et $I$ est} je dis que l'ens $\widetilde I_M$ est \uline{cofinal} dans $\widetilde I$. En effet, \deleted{soit $I'\in\widetilde I$, \deleted{\unsure{de cardinal}} \inserted{et soit} $i_0$ un majorant de $I'$ \deleted{\unsure{sur} $I$} dans $I$. Soit $I'$} $\widetilde I_M$ contient le sous-ensemble $\widetilde I_0$ formé des $I'\in\widetilde I$ qui ont un \deleted{plus} grand élément -- et $\widetilde I_0$ est bien cofinal ! Maintenant, \deleted{on peut} \unsure{appliquons|appliquer}

% archive sheet number: 102
